\documentclass[12pt, a4paper]{article}
\usepackage[UTF8]{ctex}
\usepackage{amsmath, amssymb, amsthm}
\usepackage{geometry}
\usepackage{bm}

\geometry{left=2.5cm, right=2.5cm, top=2.5cm, bottom=2.5cm}

\newcommand{\RR}{\mathbb{R}}
\newcommand{\e}{\mathrm{e}}
\newcommand{\dif}{\mathrm{d}}

\begin{document}
\section*{12.5}
\textbf{13.}证明：曲面\(z = xf\left(\dfrac{y}{x}\right)(x \neq 0)\)在任一点处的切平面都通过原点，其中\(f\)具有连续偏导数.

证明：设\(F(x, y, z) = xf\left(\dfrac{y}{x}\right) - z\)，则曲面在\(\forall (x_0, y_0, z_0)\)处的切平面方程为
\[F_x(x_0, y_0, z_0)(x - x_0) + F_y(x_0, y_0, z_0)(y - y_0) + F_z(x_0, y_0, z_0)(z - z_0) = 0.\]
即
\[\left(f\left(\frac{y_0}{x_0}\right) - \frac{y_0}{x_0}f'\left(\frac{y_0}{x_0}\right)\right)(x - x_0) + f'\left(\frac{y_0}{x_0}\right)(y - y_0) - (z - z_0) = 0.\]
将点\((0, 0, 0)\)代入上式，得
\[-x_0f\left(\frac{y_0}{x_0}\right) + y_0f'\left(\frac{y_0}{x_0}\right) - y_0f'\left(\frac{y_0}{x_0}\right) + z_0 \equiv 0.\]
故切平面过原点.

\textbf{15.}设\(F(x, y, z)\)具有连续偏导数，且\(F_x^2 + F_y^2 + F_z^2 \neq 0\).进一步，设\(k\)为正整数，\(F(x, y, z)\)为\(k\)次齐次函数，即对于任意的实数\(t\)和\((x, y, z)\)，成立
\[F(tx, ty, tz) = t^k F(x, y, z).\]
证明：曲面\(F(x, y, z) = 0\)上的所有点的切平面相交于一定点.

证明：对上式关于\(t\)求偏导，得
\[x\frac{\partial F(tx, ty, tz)}{\partial (tx)} + y\frac{\partial F(tx, ty, tz)}{\partial (ty)} + z\frac{\partial F(tx, ty, tz)}{\partial (tz)} = k t^{k-1} F(x, y, z).\]
令\(t = 1\)，则
\[xF_x(x, y, z) + yF_y(x, y, z) + zF_z(x, y, z) = k F(x, y, z).\]
由于点\((x_0, y_0, z_0)\)在曲面\(F(x, y, z) = 0\)上，故
\[x_0F_x(x_0, y_0, z_0) + y_0F_y(x_0, y_0, z_0) + z_0F_z(x_0, y_0, z_0) = 0.\]
而曲面在点\((x_0, y_0, z_0)\)处的切平面方程为
\[F_x(x_0, y_0, z_0)(x - x_0) + F_y(x_0, y_0, z_0)(y - y_0) + F_z(x_0, y_0, z_0)(z - z_0) = 0.\]
将点\((0, 0, 0)\)代入上式，得
\[-x_0F_x(x_0, y_0, z_0) - y_0F_y(x_0, y_0, z_0) - z_0F_z(x_0, y_0, z_0) \equiv 0.\]
故切平面过原点.

\section*{12.6}
\textbf{9.}求由方程\(2x^2 + 2y^2 + z^2 + 8yz - z + 8 = 0\)所确定的隐函数\(z = z(x, y)\)的极值.

解：设\(F(x, y, z) = 2x^2 + 2y^2 + z^2 + 8yz - z + 8 = 0\)，则
\[F_x = 4x, \quad F_y = 4y + 8z, \quad F_z = 2z + 8y - 1.\]
\[\frac{\partial z}{\partial x} = -\frac{F_x}{F_z} = -\frac{4x}{2z + 8y - 1}, \quad \frac{\partial z}{\partial y} = -\frac{F_y}{F_z} = -\frac{4y + 8z}{2z + 8y - 1}.\]
令\(\dfrac{\partial z}{\partial x} = 0\)且\(\dfrac{\partial z}{\partial y} = 0\)，解得\(x = 0, y = -2z\).代入原方程整理得
\[7z^2 +z - 8 = 0.\]
解得\(z = 1\)或\(z = -\dfrac{8}{7}\).故驻点为\(P_1(0, -2)\)和\(P_2\left(0, \dfrac{16}{7}\right)\).
\[\frac{\partial^2 z}{\partial x^2} = -\frac{4(2z + 8y - 1)^2 + 32x^2}{(2z + 8y - 1)^3}\]
\[\frac{\partial^2 z}{\partial x \partial y} = \frac{32x(7y - 1)}{(2z + 8y - 1)^3},\]
\[\frac{\partial^2 z}{\partial y^2} = -\frac{4(2z + 8y - 1)^2 - 16(4y + 8z)(2z + 8y - 1) + 2(4y + 8z)^2}{(2z + 8y - 1)^3}.\]
在驻点\(P_1\)处，有
\[\frac{\partial^2 z}{\partial x^2} = \frac{4}{15}, \quad \frac{\partial^2 z}{\partial x \partial y} = 0, \quad \frac{\partial^2 z}{\partial y^2} = \frac{4}{15}.\]
\[H = \frac{\partial^2 z}{\partial x^2} \cdot \frac{\partial^2 z}{\partial y^2} - \left(\frac{\partial^2 z}{\partial x \partial y}\right)^2 = \frac{16}{225} > 0.\]
故\(P_1\)为极小值点，极小值为\(z(0, -2) = 1\).
在驻点\(P_2\)处，有
\[\frac{\partial^2 z}{\partial x^2} = -\frac{4}{15}, \quad \frac{\partial^2 z}{\partial x \partial y} = 0, \quad \frac{\partial^2 z}{\partial y^2} = -\frac{4}{15}.\]
\[H = \frac{\partial^2 z}{\partial x^2} \cdot \frac{\partial^2 z}{\partial y^2} - \left(\frac{\partial^2 z}{\partial x \partial y}\right)^2 = \frac{16}{225} > 0.\]
故\(P_2\)为极大值点，极大值为\(z\left(0, \dfrac{16}{7}\right) = -\dfrac{8}{7}\).

\textbf{12.}证明：圆的所有内接\(n\)边形中，以正\(n\)边形的面积为最大.

证明：不妨设圆的半径为\(\sqrt{2}\)，内接\(n\)边形的顶点为\(P_1, \cdots P_n\).设\(\theta_i = \angle P_iOP_{i + 1}, i = 1, \cdots n - 1, \theta_n = \angle P_nOP_1\).则约束条件\(\sum\limits_{i = 1}^{n} \theta_i - 2\pi = 0\)，内接\(n\)边形的面积为
\[S = \sum_{i = 1}^{n} \sin\theta_i,\quad L = \sum_{i = 1}^{n} \sin\theta_i - \lambda\left(\sum_{i = 1}^{n} \theta_i - 2\pi\right).\]
令\(L_{\theta_i} = \cos\theta_i - \lambda = 0\)，解得\(\theta_1 = \cdots = \theta_n = \dfrac{2\pi}{n}\)，由于问题的最大值必定存在，因此它就是最大值点，即正\(n\)边形时面积最大.

\textbf{13.}证明：当\(0 < x < 1, 0 < y < +\infty\)时，成立不等式
\[yx^y(1 - x) < \e^{-1}.\]

证明：设\(f(x, y) = yx^y(1 - x)\)，则
\[f_x = y\left(yx^{y - 1}(1 - x) - x^y\right) = yx^{y - 1}(y - yx - x),\]
\[f_y = yx^y(1 - x)\ln x + x^y(1 - x).\]
令\(f_x = 0\)且\(f_y = 0\)，解得
\[x = \frac{y}{y + 1}, \quad y = -\frac{1}{\ln x}.\]
联立得\(x = \dfrac{1}{\e}, y = \e - 1\).在此点处，
\[f\left(\frac{1}{\e}, \e - 1\right) = (\e - 1)\left(\frac{1}{\e}\right)^{\e - 1}\left(1 - \frac{1}{\e}\right) = \e^{-1}.\]
又当\(x \to 0^+\)或\(x \to 1^-\)时，\(f(x, y) \to 0\)，故不等式成立.

\section*{12.7}
\textbf{1.}求下列函数的条件极值：

(1)\(f(x, y) = xy\)，约束条件为\(x + y = 1\).

解：设\(g(x, y) = x + y - 1 = 0, L = f - \lambda g.\)
\[\begin{cases}
L_x = y - \lambda = 0,\\
L_y = x - \lambda = 0,\\
g(x, y) = 0.
\end{cases}\]
解得\(x = y = \dfrac{1}{2}\)，此时\(f_{\text{max}} = f\left(\dfrac{1}{2}, \dfrac{1}{2}\right) = \dfrac{1}{4}\).

(3)\(f(x, y, z) = \dfrac{x^2}{a^2} + \dfrac{y^2}{b^2} + \dfrac{z^2}{c^2}\)，约束条件为\(\begin{cases}
x^2 + y^2 + z^2 = 1,\\
A x + B y + C z = 0.
\end{cases}\)，其中\(a > b > c > 0, A^2 + B^2 + C^2 = 1\).

解：设\(g(x, y, z) = x^2 + y^2 + z^2 - 1 = 0, h(x) = Ax + By + Cz = 0, L = f - \lambda g - \mu h.\)
\[\begin{cases}
L_x = \dfrac{2x}{a^2} - 2\lambda x - \mu A = 0,\\
L_y = \dfrac{2y}{b^2} - 2\lambda y - \mu B = 0,\\
L_z = \dfrac{2z}{c^2} - 2\lambda z - \mu C = 0,\\
g(x, y, z) = 0,\\
h(x, y, z) = 0.
\end{cases}\]
由前三个方程可得
\[f = \lambda,\]
\[\begin{cases}
x = \dfrac{\mu A a^2}{2(1 - \lambda a^2)},\\
y = \dfrac{\mu B b^2}{2(1 - \lambda b^2)},\\
z = \dfrac{\mu C c^2}{2(1 - \lambda c^2)}.
\end{cases}\]
将其代入约束条件中，解得
\[\dfrac{A^2 a^2}{1 - \lambda a^2} + \dfrac{B^2 b^2}{1 - \lambda b^2} + \dfrac{C^2 c^2}{1 - \lambda c^2} = 0.\]
即
\[\lambda^2 + \left(\frac{A^2 - 1}{a^2} + \frac{B^2 - 1}{b^2} + \frac{C^2 - 1}{c^2}\right)\lambda + \left(\frac{A^2}{b^2c^2} + \frac{B^2}{c^2a^2} + \frac{C^2}{a^2b^2}\right) = 0.\]
故\(f\)的极大极小值是方程的两个根.

\textbf{2.}在周长为\(2p\)的一切三角形中，找出面积最大的三角形.

解：设三角形的三边长为\(x, y, z\)，则约束条件为\(g(x, y, z) = x + y + z - 2p = 0\)，面积为
\[S = \sqrt{p(p - x)(p - y)(p - z)}.\]
设\(L = S - \lambda g\)，则
\[\begin{cases}
L_x = \dfrac{-p(p - y)(p - z)}{2S} - \lambda = 0,\\
L_y = \dfrac{-p(p - x)(p - z)}{2S} - \lambda = 0,\\
L_z = \dfrac{-p(p - x)(p - y)}{2S} - \lambda = 0,\\
g(x, y, z) = 0.
\end{cases}\]
由前三个方程可得\(x = y = z = \dfrac{2p}{3}\).故面积最大的三角形为正三角形，面积为
\[\frac{\sqrt{3}}{4}\left(\frac{2p}{3}\right)^2 = \frac{\sqrt{3}p^2}{9}.\]

\textbf{4.}抛物面\(z = x^2 + y^2\)被平面\(x + y + z = 1\)截成一椭圆，求原点到这个椭圆的最长距离与最短距离.

解：设椭圆上的点为\((x, y, z)\)，则约束条件为
\[g(x, y, z) = x^2 + y^2 - z = 0,\]
\[h(x, y, z) = x + y + z - 1 = 0,\]
设\(f(x, y, z) = x^2 + y^2 + z^2, L = f - \lambda g - \mu h\)，则
\[\begin{cases}
L_x = 2x - 2\lambda x - \mu = 0,\\
L_y = 2y - 2\lambda y - \mu = 0,\\
L_z = 2z + \lambda - \mu = 0,\\
g(x, y, z) = 0,\\
h(x, y, z) = 0.
\end{cases}\]
解得\((x, y, z) = \left(\dfrac{-1 + \sqrt{3}}{2}, \dfrac{-1 + \sqrt{3}}{2}, 2 - \sqrt{3}\right)\)或\((x, y, z) = \left(\dfrac{-1 - \sqrt{3}}{2}, \dfrac{-1 - \sqrt{3}}{2}, 2 + \sqrt{3}\right).\)
由于椭圆半长轴半短轴必然存在，故对应的\(b = \sqrt{f_{\text{min}}} = \sqrt{9 - 5\sqrt{3}}, a = \sqrt{f_{\text{max}}} = \sqrt{9 + 5\sqrt{3}}.\)

\textbf{7.}求平面\(Ax + By + Cz = 0\)与柱面\(\dfrac{x^2}{a^2} + \dfrac{y^2}{b^2} = 1\)相交所成的椭圆的面积（\(A, B, C\)都不为零；\(a, b\)为正数）.

解：设\(f(x, y, z) = x^2 + y^2 + z^2\)，椭圆上的点为\((x, y, z)\)，则约束条件为
\[g(x, y, z) = Ax + By + Cz = 0,\]
\[h(x, y, z) = \frac{x^2}{a^2} + \frac{y^2}{b^2} - 1 = 0,\]
设\(L = f - \lambda g - \mu h\)，则
\[\begin{cases}
\operatorname{grad} L = 0,\\
g(x, y, z) = 0,\\
h(x, y, z) = 0.
\end{cases}\]
解得\(f = \mu, C^2\mu^2 - (a^2A^2 + b^2B^2 + C^2(a^2 + b^2))\mu + a^2b^2(A^2 + B^2 + C^2) = 0\)，于是
\[S = \pi\sqrt{\mu_1}\sqrt{\mu_2} = \pi ab\frac{\sqrt{A^2 + B^2 + C^2}}{|C|}.\]

\textbf{8.}求\(z = \dfrac{1}{2}(x^4 + y^4)\)在条件\(x + y = a\)下的最小值，其中\(x \geq 0, y \geq 0, a\)为常数.并证明不等式
\[\frac{x^4 + y^4}{2} \geq \left(\frac{x + y}{2}\right)^4.\]

解：设\(f(x, y) = \dfrac{1}{2}(x^4 + y^4), g(x, y) = x + y - a, L = f - \lambda g.\)
\[\begin{cases}
\operatorname{grad} L = 0,\\
g = 0.
\end{cases}\]
解得\((x, y) = \left(\dfrac{1}{2}, \dfrac{1}{2}\right).\)而\(f_{xx}\left(\dfrac{1}{2}, \dfrac{1}{2}\right) = \dfrac{3}{2}, f_{yy}\left(\dfrac{1}{2}, \dfrac{1}{2}\right) = \dfrac{3}{2}, f_{xy}\left(\dfrac{1}{2}, \dfrac{1}{2}\right) = 0\)，故\(\left(\dfrac{1}{2}, \dfrac{1}{2}\right)\)是极小值点，极小值为\(\dfrac{a^2}{16}.\)即
\[\frac{x^4 + y^4}{2} \geq \frac{a^4}{16} = \left(\frac{x + y}{2}\right)^4.\]

\textbf{9.}当\(x > 0, y > 0, z > 0\)时，求函数
\[f(x, y, z) = \ln x + 2\ln y + 3\ln z\]
在球面\(x^2 + y^2 + z^2 = 6R^2\)上的最大值.并由此证明：当\(a, b, c\)为正实数时，成立不等式
\[ab^2c^3 \leq 108\left(\frac{a + b + c}{6}\right)^6.\]

解：设\(g(x, y, z) = x^2 + y^2 + z^2 - 6R^2, L = f - \lambda g.\)
\[\begin{cases}
\operatorname{grad} L = 0,\\
g = 0.
\end{cases}\]
解得\((x, y, z) = (R, \sqrt{2}R, \sqrt{3}R).\)故最大值\(f_{\text{max}} = 6\ln R + \ln(6\sqrt{3})\)，于是
\[\ln(xy^2z^3) \leq 6\ln R + \ln(6\sqrt{3}),\]
即
\[xy^2z^3 \leq 6\sqrt{3}R^6 = \frac{\sqrt{3}}{36}(x^2 + y^2 + z^2)^3.\]
令\(a = x^2, b = y^2, c = z^2\)，得到
\[ab^2c^3 \leq 108\left(\frac{a + b + c}{6}\right)^6.\]

\textbf{10.}(1)求函数\(f(x, y, z) = x^a y^b, z^c (x > 0, y > 0, z > 0)\)在约束条件\(x^k + y^k + z^k = 1\)下的极大值，其中\(k, a, b, c\)均为正常数.

(2)利用(1)的结果证明：对于任何正数\(u, v, w\)，成立不等式
\[\left(\frac{u}{a}\right)^a \left(\frac{v}{b}\right)^b \left(\frac{w}{c}\right)^c \leq \left(\frac{u + v + w}{a + b + c}\right)^{a + b + c}.\]

解：(1)设\(g(x, y, z) = x^k + y^k + z^k - 1, L = f - \lambda g.\)
\[\begin{cases}
\operatorname{grad} L = 0,\\
g = 0.
\end{cases}\]
解得\((x, y, z) = \left(\left(\dfrac{a}{a + b + c}\right)^\frac{1}{k}, \left(\dfrac{b}{a + b + c}\right)^\frac{1}{k}, \left(\dfrac{c}{a + b + c}\right)^\frac{1}{k}\right).\)代入得
\[f_{\text{max}} = \left(\frac{a^a b^b c^c}{(a + b + c)^{a + b + c}}\right)^\frac{1}{k}.\]

(2)取\(k = 1\)，得到
\[x^a y^b z^c \leq \frac{a^a b^b c^c}{(a + b + c)^{a + b + c}}.\]
令\((x, y, z) = \left(\left(\dfrac{u}{u + v + w}\right)^a, \left(\dfrac{v}{u + v + w}\right)^b, \left(\dfrac{w}{u + v + w}\right)^c\right)\)，代入得
\[\left(\dfrac{u}{u + v + w}\right)^a \left(\dfrac{v}{u + v + w}\right)^b \left(\dfrac{w}{u + v + w}\right)^c \leq \frac{a^a b^b c^c}{(a + b + c)^{a + b + c}},\]
即
\[\left(\frac{u}{a}\right)^a \left(\frac{v}{b}\right)^b \left(\frac{w}{c}\right)^c \leq \left(\frac{u + v + w}{a + b + c}\right)^{a + b + c}.\]

\textbf{13.}设\(a_1, a_2, \cdots, a_n\)为\(n\)个已知正数.求\(n\)元函数
\[f(x_1, x_2, \cdots, x_n) = \sum_{k = 1}^{n} a_k x_k\]
在约束条件
\[\sum_{k = 1}^{n} x_k^2 \leq 1\]
下的最大值和最小值.

解：由Cauchy-Schwarz不等式，成立
\[\left(\sum_{k = 1}^{n} a_k x_k\right)^2 \leq \left(\sum_{k = 1}^{n} a_k^2\right)\left(\sum_{k = 1}^{n} x_k^2\right) \leq \left(\sum_{k = 1}^{n} a_k^2\right).\]
即
\[-\sqrt{\sum_{k = 1}^{n} a_k^2} \leq \sum_{k = 1}^{n} a_k x_k \leq \sqrt{\sum_{k = 1}^{n} a_k^2}.\]
当\(x_k = \dfrac{a_k}{\sqrt{a_k^2}}\)时，\(\sum\limits_{k = 1}^{n} x_k^2 = 1, \left(\sum\limits_{k = 1}^{n} a_k x_k\right)_{\text{max}} = \sqrt{\sum\limits_{k = 1}^{n} a_k^2}\)；当\(x_k = -\dfrac{a_k}{\sqrt{a_k^2}}\)时，\(\sum\limits_{k = 1}^{n} x_k^2 = 1, \left(\sum\limits_{k = 1}^{n} a_k x_k\right)_{\text{min}} = -\sqrt{\sum\limits_{k = 1}^{n} a_k^2}\).

\textbf{14.}求二次型\(\sum_{i, j = 1}^{n} a_{ij}x_ix_j(a_{ij} = a_{ji})\)在\(n\)维单位球面
\[\left\{(x_1, x_2, \cdots, x_n) \in \RR^n \left|\sum_{k = 1}^{n} x_k^2 = 1\right.\right\}\]
上的最大值与最小值.

解：设\(A = (a_{ij}), X = (x_1, x_2, \cdots, x_n)^\mathrm{T}, \sum_{i, j = 1}^{n} a_{ij}x_ix_j = X^\mathrm{T}AX = f(X)\)，约束条件\(X^\mathrm{T} = 1\).

由\(A\)是实对称矩阵，存在正交矩阵\(P\)满足
\[P^\mathrm{T}AP = \Lambda = \operatorname{diag}(\lambda_1, \lambda_2, \cdots, \lambda_n),\]
其中\(\lambda_1 \leq \lambda_2 \leq \cdots \leq \lambda_n\)是\(A\)的特征值.
令\(Y = P^\mathrm{T}X, X = PY,\)且
\[X^\mathrm{T}X = Y^\mathrm{T}P^\mathrm{T}PY = Y^\mathrm{T}Y = 1,\]
\[F(X) = X^\mathrm{T}AX = Y^\mathrm{T}P^\mathrm{T}APY = Y^\mathrm{T}\Lambda Y = \sum_{i = 1}^{n} \lambda_i y_i^2.\]
在约束\(Y^\mathrm{T}Y = \sum_{i = 1}^{n} y_i^2 = 1\)下，有
\[\lambda_1 \leq  \sum_{i = 1}^{n} \lambda_i y_i^2 \leq \lambda_n.\]
当且仅当\(Y = (1, 0, \cdots, 0)^\mathrm{T}\)时取最小值\(\lambda_1\)；\(Y = (0, \cdots, 0, 1)^\mathrm{T}\)时取最大值\(\lambda_n\).

\end{document}